% ECONOMICS_PROBLEM: {"id": "D91-482", "title": "Economic models of costly computation", "jel_code": "D91", "source_id": "OP-0275"}
% !TeX program = xelatex
\documentclass[11pt,a4paper]{article}
\usepackage[margin=23mm]{geometry}
\usepackage{fontspec}
\setmainfont{Latin Modern Roman}
\usepackage{amsmath,amssymb,mathtools}
\usepackage{xcolor,enumitem,booktabs,longtable,array,xurl,hyperref}
\definecolor{muted}{HTML}{53636C}
\hypersetup{colorlinks=true,urlcolor=blue,linkcolor=blue}
\setlength{\parindent}{0pt}
\setlength{\parskip}{5pt}
\newcommand{\R}{\mathbb R}
\newcommand{\E}{\mathbb E}
\newcommand{\N}{\mathbb N}
\newcommand{\ind}{\mathbf 1}
\newcommand{\Var}{\operatorname{Var}}
\newcommand{\Cov}{\operatorname{Cov}}
\newcommand{\supp}{\operatorname{supp}}
\DeclareMathOperator*{\argmin}{arg\,min}
\DeclareMathOperator*{\argmax}{arg\,max}
\newcommand{\entrymeta}[3]{{\sffamily\small\color{muted}\textbf{\texttt{#1}}\quad|\quad #2\par #3\par}}
\newcommand{\Problemref}[1]{\texttt{#1}}
\newcommand{\JELref}[2]{\texttt{#2} (JEL \texttt{#1})}
\begin{document}
\section*{Economic models of costly computation}
\textbf{Problem ID:} \texttt{D91-482}\quad \textbf{JEL:} \texttt{D91}\par
% BEGIN ECONOMICS STATEMENT
\label{problem:OP-0275}
{\sffamily\small\color{muted}Original subject group: Economic theory, equilibrium and social choice\par}
\label{definitions:problem:30007556}
\textbf{Audit classification:} Published research agenda. Verified computational-constraint directions in the survey’s §7. The cost-and-repertoire identification exercise is editorial; finite maximization is elementary.
\subsection*{Definitions and precise research target}
\textbf{Model and research target.} Consider economic decision problems \(d\) drawn from a distribution \(\mathcal D\). Each problem specifies a finite feasible set \(X_d\), payoff function \(u_d:X_d\to[0,1]\), and an information representation accessible through queries. A decision algorithm \(A\) maps query histories into further queries or a final choice \(x\). Let \(C(A,d)\) be its expected number of elementary operations, and let \(\kappa\geq0\) be an individual's cost per operation. A computational-choice model selects an algorithm from a stated repertoire \(\mathcal A\) to maximize
\[
 V(A,d,\kappa)=E[u_d(A(d))]-\kappa C(A,d).
\]
The target is to construct and identify a parsimonious repertoire and cost specification that jointly predict choices, response times, and the effects of experimentally varied information complexity. Estimate \(\kappa\) and repertoire heterogeneity using randomized changes in representation while holding material payoffs fixed. Determine restrictions under which algorithmic cost can be distinguished from preferences over outcomes and statistical choice noise. Require successful prediction on new problem families and an explicit account of how primitive computational operations map into measured time. A solution may establish failure of portable parameterization and identify the additional primitives required.

\emph{Definitions and maintained assumptions (editorial unless a source definition is named).} Fix an explicit machine model, query alphabet, encoded problem family \(\mathcal D\), and repertoire of algorithms. An algorithm must terminate almost surely; its random output is feasible and its expected operation count is finite. Costs and payoffs are expressed in comparable units, and the map from operation count to response time is part of the model. A representation treatment changes the encoding or query access while preserving outcomes and objective payoffs; this does not by itself prove invariance of utility or algorithm availability. To avoid assuming the decision maker sees the full problem before selecting a procedure, either let a procedure be chosen ex ante to maximize \(\mathbb E_{d\sim\mathcal D}V(A,d,\kappa)\), or state the information available at the selection node. The unknown parameters include repertoire membership, \(\kappa\), utility curvature if not fixed, and time measurement error. Identification is uniqueness, or a sharp set, among these parameters generating the same action-time law across treatments. With a known finite repertoire, optimizing \(V\) by enumeration already solves the purely numerical choice problem.
\subsection*{Variants and assumption changes}
\begin{enumerate}
\item \textbf{Fixed repertoire (editorial).} Specify finitely many query procedures with experimentally calibrated time costs. Identify cost heterogeneity from representation changes and test predictions on new distributions of encoded problems.
\item \textbf{Uniform complexity (editorial).} Let input size grow and restrict the repertoire to programs of bounded description length or operation budget. Study payoff loss as a function of that budget and whether parameters estimated at one size predict choice at another.
\end{enumerate}
\subsection*{References and source audit}
\begin{enumerate}
\item §7 pp.53–54 in February 2023 manuscript. Supports the stated research area; exact displayed specification is editorial. DOI publication is2024; full inspected manuscript is February2023. \url{https://geoffroydeclippel.net/Working%20Papers%20PDFs/JEL.pdf}
\end{enumerate}
\begin{enumerate}\raggedright
\item\hypertarget{definitions:source:30007556:1}{} Geoffroy de Clippel; Kareen Rozen. 2024. \emph{Bounded Rationality in Choice Theory: A Survey}. Published JEL 62(3),995–1039 (2024); inspected February2023 author manuscript §§6.4–7, printed pp.51–54.
\url{https://geoffroydeclippel.net/Working%20Papers%20PDFs/JEL.pdf}.
DOI: \url{https://doi.org/10.1257/jel.20231592}.
\end{enumerate}

\subsection*{Common conventions: Source mathematical conventions}

These conventions apply to each entry unless explicitly replaced.

\textbf{Models and identification.} A primitive model \(m\) specifies all probability laws, feasible choices, information, equilibrium and measurement rules used by its target. The admissible class \(\mathcal M\) is fixed independently of outcomes. If \(P_O(m)\) is its observable probability law and \(\tau(m)\) the target, its identified set at observable law \(P\) is
\[
 \mathcal I_\tau(P)=\{\tau(m):m\in\mathcal M,\ P_O(m)=P\}.
\]
Point identification means that this set is a singleton. An empty set signals incompatibility of data and assumptions. Coordinate bounds need not describe the joint set. Bounds are sharp only if every retained target is attainable or arbitrarily approachable by compatible models. Finite bounds require explicit restrictions on unobserved quantities; unrestricted confounding, hidden wealth, extrapolation or arbitrary equilibrium selection can yield uninformative sets.

\textbf{Causal interventions.} A potential outcome \(Y_i(p)\) is the outcome under a complete policy regime \(p\), including timing, financing, information and market spillovers. The target population and units are fixed before treatment. With interference, use the complete assignment vector or a justified exposure mapping; individual no-interference notation alone cannot identify an economy-wide intervention. A difference of mean potential outcomes does not require the cross-world joint distribution. Conditioning on post-treatment migration, survival, enrollment or employment changes the estimand and needs separate justification.

\textbf{Statistical statements.} An estimator, confidence set, forecast or test specifies a sampling law, model class and loss/coverage criterion. Uniform \((1-\alpha)\)-coverage means \(\inf_{m\in\mathcal M}\Pr_m\{\tau(m)\in C_n\}\geq1-\alpha\), or an explicitly asymptotic version. The whole parameter space trivially has coverage, so informativeness requires a stated width, risk or power objective. A forecasting improvement is relative to a named benchmark, information set and evaluation population.

\textbf{Equilibrium and optimization.} An equilibrium comprises feasible plans, optimality, beliefs where needed, budget constraints and all market/resource clearing. The primitives or a stated benchmark define each correspondence \(\mathcal E(p;m)\). Nonemptiness is assumed or proved, never inferred just from compact policy sets. A selected-equilibrium target uses a declared selection rule; a robust target quantifies over all permitted equilibria. Use suprema or infima unless attainment is established. An \(\varepsilon\)-optimal policy is feasible and within \(\varepsilon\) of the finite optimum value. Report an empty feasible set separately.

\textbf{Welfare and accounting.} Normative utility, interpersonal normalization, social weights, numeraire, discounting and the use of public receipts are primitives. Behavioral utility need not coincide with the welfare criterion. Household welfare includes ownership income and rents once; adding profits again to compensating variation generally double-counts them. A monetary equivalent is defined by an explicit transfer and price convention, with monotonicity and range conditions. Average effects, mechanism contributions, transfers and welfare losses are different targets. Infinite sums require convergence and dynamic feasible plans require terminal or no-Ponzi conditions.

\textbf{Source versions.} A working-paper date, revision date and journal publication year are distinguished. A DOI for a journal version is not automatically the DOI of an earlier manuscript. Locators refer to the stated version and printed pages where specified. A metadata-only check does not verify a claim about a full-text conclusion. Every entry's source subsection narrows the provenance claim accordingly.
% END ECONOMICS STATEMENT
\end{document}
